\section{Introduction}\label{sec:introduction}
Financial and economic decisions depend on how people and organizations
respond to one another under uncertainty. A protocol designer considering an
incentive, for example, may need to understand not only its direct cost but
also whether partners participate, how other participants react, and how
those reactions affect later decisions. Comparing such choices requires
scenarios that make actors, constraints, and assumed events explicit while
showing how their interactions produce different outcomes.

Agent-based modeling studies a system through individual participants and
their interaction rules; its conclusions depend on both the model assumptions
and their implementation \cite{grimm2020,galan2009}. Language-model agents can
condition choices on context \cite{park2023,yang2025}, but an expressed choice
alone does not determine its economic consequences. If an actor discloses
information, the model must still specify who receives it, when they may react,
and what changes in resources or state follow. These distinctions are needed
to interpret why a simulated path occurred and what assumptions support it.

We present Boyun, a framework for exploring financial and economic decisions
through explicit simulated worlds. A world specifies participants, their
resources and information, available actions, and rules governing action
consequences. Participants receive restricted observations and a decision
model selects from feasible actions; an execution engine applies the specified
rules and records the resulting transitions. This separation lets an analyst
inspect the supplied information, selected actions, and computed outcomes,
and replay execution using recorded choices. It does not expose the decision
model's internal deliberation or establish that its choices represent real
people.

The paper asks: \emph{Can an explicitly specified multi-agent world constrain
information, spending, and settlement while making the effects of alternative
information and interaction rules inspectable?} Our initial scope includes
capital allocation, participation incentives, information policies, and stress
assumptions in crypto and decentralized finance. The intended audience is
researchers, protocol designers, and risk analysts comparing conditional
outcomes under declared assumptions.

\begin{figure*}[t]
\centering
\begin{tikzpicture}[node distance=7mm and 8mm]
\node[box] (build) {Alice builds in round one\\Declared cost: 20 cents};
\node[box,right=of build] (message) {A subscribed message\\reaches Bob in round two};
\node[box,right=of message] (join) {Bob joins after delivery\\Declared cost: 30 cents};
\draw[arr] (build)--(message);\draw[arr] (message)--(join);
\end{tikzpicture}
\caption{A controlled example of mechanism attribution. With the same
programmed participant policies, readiness is 6 points with message delivery
and a joint-action rule, 3 with delivery alone, and 1 without delivery.
Section~\ref{sec:ablation} defines the units and all four conditions. The
responses are imposed test policies, not observations of people.}
\label{fig:teaser}
\end{figure*}

We investigate this question through six synthetic experiments, including a
controlled example in Fig.~\ref{fig:teaser} that separates the effects of
message delivery and an interaction rule. The experiments check calculations
and execution under specified participant policies; they do not establish
human behavioral fidelity, market forecasting accuracy, or improved user
decisions. Those stronger claims require separate empirical evaluation against
appropriate baselines.

The paper contributes (1) a world and observation specification connecting
participant goals, information, feasible actions, and explicit consequences;
(2) a synchronous execution procedure with recorded decisions and exact replay;
and (3) synthetic experiments that test the specified arithmetic and
interaction mechanisms. An optional NEAR integration supports artifact
registration; it does not validate the behavior represented by the simulation.

\section{Related Work}\label{sec:related}
\subsection{Agent-Based Modeling and Model Description}
The Overview, Design concepts, and Details (ODD) protocol is a structured
method for describing simulation models \cite{grimm2020}. It asks authors to
state their entities, initialization, schedules, input data, and submodels so
that others can understand and reproduce them. Boyun follows this principle
by requiring an explicit world specification before execution. Gal\'an
et al. distinguish modeling errors and implementation artifacts in agent-based
systems \cite{galan2009}; this motivates separate checks of numerical behavior
and the realism of the assumptions that behavior implements.

\subsection{Economic Environments and Heterogeneous Decisions}
Economic simulation connects decisions to an explicit allocation mechanism,
not merely to a conversation between roles. The AI Economist combines an
economic environment with two levels of reinforcement learning: individual
agents adapt their economic behavior while a social planner learns tax policy
\cite{zheng2020}. Its Foundation framework separates agents, environment
components, and scenarios. The relevant principle is that an individual's
utility and the policy designer's objective are different objects. Boyun
similarly declares each participant's preferences separately from the user's
research question, but does not train policies or optimize tax schedules.
The AI Economist's experimental results do not transfer to Boyun.

EconAgent uses language-model agents to make work and consumption decisions
inside a macroeconomic environment with taxation, goods, wages, savings, and
interest-rate rules \cite{li2024econagent}. Perception and memory condition
heterogeneous decisions on personal and economic context. Its environment
then computes economic consequences. This is a close precedent for separating
an agent decision mechanism from an economic transition model. Boyun instead
uses reviewed finite action menus, explicit utility weights, and native typed
Jev answers for bounded financial and economic questions. It does not implement
EconAgent's labor/consumption economy or inherit its reported stylized facts.
EconAgent itself distinguishes reproducing such regularities from accurate
forecasting; that distinction also applies here.

Market microstructure requires a more specific execution environment. ABIDES,
the Agent-Based Interactive Discrete Event Simulation framework, represents
trading agents, an exchange, protocol messages, and configurable network
latency \cite{byrd2019}. It illustrates what must be specified to study order
execution or market impact. Boyun's strategic interface does not contain that
exchange mechanism. Analysis of Uniswap describes automated market makers
(AMMs), whose reserve rules determine exchange prices \cite{angeris2021};
this informs the separate financial laboratory, not every strategic world.
A hypothetical participation score is neither a transaction nor a return.

\begin{table*}[t]
\caption{Relevant research systems and their relationship to Boyun. These are
method comparisons, not measured rankings. None of these additional projects
is a dependency of Boyun's active strategic engine. MiroFish remains the direct
software provenance described separately.}
\label{tab:economic-related}
\centering\footnotesize
\begin{tabularx}{\textwidth}{@{}p{0.15\textwidth}p{0.28\textwidth}YY@{}}
\toprule System & Explicit mechanism & Relevant design or evaluation principle & Boundary in Boyun \\
\midrule
AI Economist / Foundation \cite{zheng2020} & Economic environment with individual agents and a policy-setting planner & Separate actor utility from a designer's aggregate objective & No learned planner or two-level policy training \\
EconAgent \cite{li2024econagent} & Work/consumption choices within numeric macroeconomic rules & Heterogeneity needs economic context and an environment that settles choices & No integrated macroeconomy or inherited validation \\
ABIDES \cite{byrd2019} & Message-based exchange simulation and latency & Execution rules are necessary for price, fill, and impact claims & No limit-order-book exchange in the strategic interface \\
OpenSpiel \cite{lanctot2019} & Explicit game states, legal actions, information and returns & Evaluate policy behavior using known games and compatible solution metrics & No imported solvers or large dynamic-game solution \\
Repeated-game experiments \cite{akata2025} & Language models choose in controlled payoff games & Test coordination and response to opponent information; do not assume rationality & Results for other models do not establish Jev behavior \\
\bottomrule
\end{tabularx}
\end{table*}

\subsection{Game Structure and Behavioral Evaluation}
A finite game specifies participants, available strategies, and payoffs.
A Nash equilibrium is a strategy profile in which no participant benefits
from changing its own strategy alone \cite{nash1950}. Polak's lectures explain
best responses, the removal of strictly inferior strategies, and coordination
failures using small games \cite{polak3,polak5}. Boyun uses these concepts to
inspect supplied incentives; it does not assume that a model's selected
action is an equilibrium.

OpenSpiel provides game environments and algorithms for reinforcement
learning, search, and planning, with legal actions, information states,
transitions, and returns \cite{lanctot2019}. Its treatment of simultaneous,
sequential, and imperfect-information games is relevant to specifying exactly
which game a policy is being tested on. Boyun implements a narrower world
contract and small finite-game diagnostics rather than wrapping OpenSpiel's
algorithms. Evaluating Jev on small, known-payoff games before expanding the
cast would make violations and coordination failures measurable; this is a
proposed comparison, not an experiment completed here.

Akata et al. study language models in finitely repeated games and report
different behavior in self-interested and coordination settings
\cite{akata2025}. Their experiments also examine sensitivity to opponent
information. Such findings motivate testing a persona policy's behavior
instead of treating an instruction to maximize utility as proof that it does
so. Boyun's Jev interface supplies action scores without a reasoning trace;
it does not reproduce conversational reasoning interventions from that study.
Observed choices, payoff comparisons, and controlled interventions are the
available evaluation objects.

Bayesian updating revises a prior belief using the likelihood of new evidence
under each hypothesis \cite{pishronik2014}. It provides a disciplined way to
update declared assumptions, but cannot supply missing likelihoods. Proper
scoring rules evaluate issued probabilities against resolved outcomes: the
Brier score measures squared probability error, while logarithmic loss
penalizes confident errors more strongly \cite{brier1950,gneiting2007}.
These concepts support evaluation, rather than turning model confidence into
an established forecast probability.

\subsection{Contextual Agents and the Social-Simulation Precedent}
The Transformer introduced by Vaswani et al. uses attention to combine token
representations according to contextual compatibility \cite{vaswani2017}.
Co-authored by Illia Polosukhin, this work is a foundation for many modern
language models; it is not a study of agent-based forecasting. Neural
attention operates inside a model. Boyun's information graph instead controls
which messages are supplied to a participant before inference. Keeping these
levels distinct prevents a token-weighting mechanism from being mistaken
for an empirically identified model of social influence. Boyun treats its
model backend as replaceable and does not assume undocumented Jev internals.

Generative Agents combines memory, reflection, and planning to produce
context-sensitive behavior in a 25-agent social sandbox \cite{park2023}.
It provides an architectural precedent for conditioning decisions on personal
experience. Boyun uses bounded history and delivered messages, without
implementing that work's full memory and reflection architecture.

Open Agent Social Interaction Simulations (OASIS) provides simulated social
platforms with participant actions, recommendation mechanisms, and evolving
interactions \cite{yang2025}. Its large-scale experiments illustrate why the
environment and information-delivery process matter, rather than only the
language model. CAMEL, Communicative Agents for ``Mind'' Exploration of Large
Language Model Society, studies role-based cooperation between conversational
agents \cite{camel2023}. Its role-playing framework and associated software
support agent coordination; neither CAMEL nor OASIS establishes the validity
of Boyun's financial assumptions. Boyun does not execute these frameworks in
its active simulation engine.

Self-report-grounded agent research constructs agents from interviews and
survey responses and compares them with held-out responses from 1,052
participants \cite{park2026}. Its relevance is the combination of grounding
data and a defined validation target. Boyun currently lacks comparable data
for crypto participants. A detailed fictional persona, or an observed wallet
trade, does not by itself reveal a person's motivation.

\subsection{MiroFish and Software Provenance}\label{sec:mirofish}
MiroFish is the direct software and product precedent for Boyun
\cite{mirofish}. Its workflow transforms seed materials into a graph of
entities and relationships, generates participant profiles, executes a social
simulation, and assembles a report. Zep is the memory and knowledge-graph
service used for the graph stage. OASIS supplies the social environment and
CAMEL-related software supplies agent infrastructure. The OASIS software is
explicitly credited by MiroFish \cite{oasiscode}.

\begin{figure*}[t]
\centering
\begin{tikzpicture}[node distance=6mm and 8mm]
\node[box] (seed) {MiroFish seed material\\Documents and question};
\node[box,right=of seed] (graph) {Zep graph\\Entities and relationships};
\node[box,right=of graph] (cast) {Generated participants\\and environment settings};
\node[box,below=of cast] (social) {OASIS social environment\\with agent actions};
\node[box,left=of social] (history) {Interaction history\\and memory updates};
\node[box,left=of history] (report) {Report agent\\Analysis and interviews};
\draw[arr] (seed)--(graph);\draw[arr] (graph)--(cast);
\draw[arr] (cast)--(social);\draw[arr] (social)--(history);\draw[arr] (history)--(report);
\end{tikzpicture}
\caption{MiroFish's construction, social execution, and reporting workflow,
reconstructed from its published software \cite{mirofish}. This figure describes
the upstream system; Fig.~\ref{fig:architecture} describes Boyun.}
\label{fig:mirofish}
\end{figure*}

Boyun is developed in an AGPL-3.0 MiroFish fork and retains the upstream
license and attribution. The inherited browser frontend, backend, and
graph/social integrations are distinct from the active Boyun engine and, as of
the 2026-10-31 release package, are no longer distributed with it; they remain
readable in the pinned upstream commit. The
design takes the construction--participants--environment--report workflow
as a starting point, then makes action menus, numeric consequences, and replay
explicit. It does not import OASIS recommendation algorithms, Zep graph memory,
or their reported experimental results into the active execution path.

\begin{table}[t]
\caption{Execution boundaries; this is not a performance benchmark.}
\label{tab:comparison}
\centering\footnotesize
\begin{tabularx}{\columnwidth}{@{}p{0.18\columnwidth}YY@{}}
\toprule & MiroFish & Boyun \\
\midrule
Construction & Seed materials, graph, profiles & Reviewed participants and rules \\
Environment & Social-platform simulation & Synchronous numeric state \\
Decision & Platform actions & Typed affordable choices \\
Analysis & Report agent and interviews & Recorded paths, costs, replay \\
Memory & Graph integration & Bounded history and subscriptions \\
\bottomrule
\end{tabularx}
\end{table}

\subsection{Decentralized Execution and Service Guarantees}
NEAR's Nightshade design partitions blockchain state and execution among
shards while organizing their work in one logical chain \cite{nightshade2024}.
Its distinction between state validity and data availability is relevant to
an experiment registry: registering an artifact fingerprint does not make
the underlying experiment available or correct. Boyun uses NEAR as a registry,
without implementing Nightshade's consensus or sharding in its off-chain engine.

Decentralized Confidential Machine Learning (DCML) describes a NEAR AI design
for inference using trusted execution environments, attested connections,
and model-identity checks \cite{dcml}. Proof of Response addresses a different
property: obtaining a response, evidence of a broken communication edge, or
delay compensation within its protocol assumptions \cite{proofresponse2025}.
These designs distinguish confidentiality, identity, availability, and content
validity. They inform Boyun's trust-boundary analysis and future deployment
requirements; their guarantees are not established by a normal model API call
or a local timeout.
